% Options for packages loaded elsewhere
% Options for packages loaded elsewhere
\PassOptionsToPackage{unicode}{hyperref}
\PassOptionsToPackage{hyphens}{url}
\PassOptionsToPackage{dvipsnames,svgnames,x11names}{xcolor}
%
\documentclass[
  letterpaper,
  DIV=11,
  numbers=noendperiod]{scrartcl}
\usepackage{xcolor}
\usepackage{amsmath,amssymb}
\setcounter{secnumdepth}{-\maxdimen} % remove section numbering
\usepackage{iftex}
\ifPDFTeX
  \usepackage[T1]{fontenc}
  \usepackage[utf8]{inputenc}
  \usepackage{textcomp} % provide euro and other symbols
\else % if luatex or xetex
  \usepackage{unicode-math} % this also loads fontspec
  \defaultfontfeatures{Scale=MatchLowercase}
  \defaultfontfeatures[\rmfamily]{Ligatures=TeX,Scale=1}
\fi
\usepackage{lmodern}
\ifPDFTeX\else
  % xetex/luatex font selection
\fi
% Use upquote if available, for straight quotes in verbatim environments
\IfFileExists{upquote.sty}{\usepackage{upquote}}{}
\IfFileExists{microtype.sty}{% use microtype if available
  \usepackage[]{microtype}
  \UseMicrotypeSet[protrusion]{basicmath} % disable protrusion for tt fonts
}{}
\makeatletter
\@ifundefined{KOMAClassName}{% if non-KOMA class
  \IfFileExists{parskip.sty}{%
    \usepackage{parskip}
  }{% else
    \setlength{\parindent}{0pt}
    \setlength{\parskip}{6pt plus 2pt minus 1pt}}
}{% if KOMA class
  \KOMAoptions{parskip=half}}
\makeatother
% Make \paragraph and \subparagraph free-standing
\makeatletter
\ifx\paragraph\undefined\else
  \let\oldparagraph\paragraph
  \renewcommand{\paragraph}{
    \@ifstar
      \xxxParagraphStar
      \xxxParagraphNoStar
  }
  \newcommand{\xxxParagraphStar}[1]{\oldparagraph*{#1}\mbox{}}
  \newcommand{\xxxParagraphNoStar}[1]{\oldparagraph{#1}\mbox{}}
\fi
\ifx\subparagraph\undefined\else
  \let\oldsubparagraph\subparagraph
  \renewcommand{\subparagraph}{
    \@ifstar
      \xxxSubParagraphStar
      \xxxSubParagraphNoStar
  }
  \newcommand{\xxxSubParagraphStar}[1]{\oldsubparagraph*{#1}\mbox{}}
  \newcommand{\xxxSubParagraphNoStar}[1]{\oldsubparagraph{#1}\mbox{}}
\fi
\makeatother


\usepackage{longtable,booktabs,array}
\usepackage{calc} % for calculating minipage widths
% Correct order of tables after \paragraph or \subparagraph
\usepackage{etoolbox}
\makeatletter
\patchcmd\longtable{\par}{\if@noskipsec\mbox{}\fi\par}{}{}
\makeatother
% Allow footnotes in longtable head/foot
\IfFileExists{footnotehyper.sty}{\usepackage{footnotehyper}}{\usepackage{footnote}}
\makesavenoteenv{longtable}
\usepackage{graphicx}
\makeatletter
\newsavebox\pandoc@box
\newcommand*\pandocbounded[1]{% scales image to fit in text height/width
  \sbox\pandoc@box{#1}%
  \Gscale@div\@tempa{\textheight}{\dimexpr\ht\pandoc@box+\dp\pandoc@box\relax}%
  \Gscale@div\@tempb{\linewidth}{\wd\pandoc@box}%
  \ifdim\@tempb\p@<\@tempa\p@\let\@tempa\@tempb\fi% select the smaller of both
  \ifdim\@tempa\p@<\p@\scalebox{\@tempa}{\usebox\pandoc@box}%
  \else\usebox{\pandoc@box}%
  \fi%
}
% Set default figure placement to htbp
\def\fps@figure{htbp}
\makeatother





\setlength{\emergencystretch}{3em} % prevent overfull lines

\providecommand{\tightlist}{%
  \setlength{\itemsep}{0pt}\setlength{\parskip}{0pt}}



 


\usepackage{hyperref} \usepackage{amsmath}
\KOMAoption{captions}{tableheading}
\makeatletter
\@ifpackageloaded{caption}{}{\usepackage{caption}}
\AtBeginDocument{%
\ifdefined\contentsname
  \renewcommand*\contentsname{Table of contents}
\else
  \newcommand\contentsname{Table of contents}
\fi
\ifdefined\listfigurename
  \renewcommand*\listfigurename{List of Figures}
\else
  \newcommand\listfigurename{List of Figures}
\fi
\ifdefined\listtablename
  \renewcommand*\listtablename{List of Tables}
\else
  \newcommand\listtablename{List of Tables}
\fi
\ifdefined\figurename
  \renewcommand*\figurename{Figure}
\else
  \newcommand\figurename{Figure}
\fi
\ifdefined\tablename
  \renewcommand*\tablename{Table}
\else
  \newcommand\tablename{Table}
\fi
}
\@ifpackageloaded{float}{}{\usepackage{float}}
\floatstyle{ruled}
\@ifundefined{c@chapter}{\newfloat{codelisting}{h}{lop}}{\newfloat{codelisting}{h}{lop}[chapter]}
\floatname{codelisting}{Listing}
\newcommand*\listoflistings{\listof{codelisting}{List of Listings}}
\makeatother
\makeatletter
\makeatother
\makeatletter
\@ifpackageloaded{caption}{}{\usepackage{caption}}
\@ifpackageloaded{subcaption}{}{\usepackage{subcaption}}
\makeatother
\usepackage{bookmark}
\IfFileExists{xurl.sty}{\usepackage{xurl}}{} % add URL line breaks if available
\urlstyle{same}
\hypersetup{
  colorlinks=true,
  linkcolor={blue},
  filecolor={Maroon},
  citecolor={Blue},
  urlcolor={Blue},
  pdfcreator={LaTeX via pandoc}}


\author{}
\date{}
\begin{document}


\section{Foreword}\label{foreword}

This text doc is intended to be a general info/log place prior to
dissertation start.

\section{Dissertation information:}\label{dissertation-information}

Supervisor: Dr Francesco Turci Co-supervisor: Prof Stephen Clark Lab
partner: Hannah Sobieralska Project: Driving non-markovianity in open
systems Code: TP-M09 Description:

An open system exchanges matter and energy with the environment.
Typically, one distinguishes the system from the environment by
separating scales: in space, the system is much smaller than the
environment, and in time, the environment relaxes much more quickly than
the system (it Markovian, i.e.~memoryless). The two are, however, always
coupled, with interfaces whose contributions are normally taken to be
small. However, in open quantum systems, a clear separation between the
system and the environment becomes more challenging, both conceptually
and practically. Conceptually, the coupling between the system and the
bath may be strong enough to lead to important memory effects, where the
evolution of the system is strongly correlated in time, and the
environment itself behaves differently from a purely thermalised state;
from the practical point of view, computer modelling of open systems
typically requires explicitly modelling both, i.e.~simulating a closed
system and examining two subsystems (the original system of interest and
the remainder, intrepreted as the environment) {[}1{]}. In this project,
we will investigate how such memory effects develop in computational
models of open quantum systems, exploiting the Matrix Product State
(MPS) formalism and tensor network methods {[}2{]}. We will explore how
different models of system-bath coupling and structures affect the
development of non-Markovian correlations both in the system and in the
environment {[}3{]}. We will start by considering equilibration towards
the thermal state and then focus on periodically driven couplings, which
introduce additional timescales and allow us to observe the system
stroboscopically (Floquet setting). This will take us to tackle
situations in which dissipation coupled to periodic driving can lead to
the emergence of non-trivial states, such as time crystals. You will
build your project on in-house code in the dedicated, high-performance,
scientific programming language Julia, along with existing advanced
packages for the simulation of open quantum systems using tensor
networks (ITensors). You will have the opportunity to explore both
analytically tractable toy models and more complex simulation settings.
You will have the opportunity to work on this frontier topic with
support from David Strachan, a recently appointed EPSRC Postdoctoral
Pathway Fellow working in this area. For further discussion of the
project scope, please contact the project supervisor.

\begin{enumerate}
\def\labelenumi{\arabic{enumi}.}
\tightlist
\item
  Breuer and Petruccione, F. (2002). The Theory of Open Quantum Systems.
  URL 2. Brenes, et al.~Physical Review X 10, 031040, 2020URL
\item
  Strachan et al.~Physical Review Letters 134, 220403, 2025 URL
\end{enumerate}

\section{Log:}\label{log}

\subsection{27/04/2026:}\label{section}

Have received project assignments, am quite happy since this was my
second choice, has amazing supervisors, and lab partner allocation is
good too; Hannah's nice, hard working (and intelligent; she was able to
do the LFU stuff whilst ill and not having done particle physics) and i
know her. Besides ik she's taking complex disordered matter which'll be
a boon.

Next steps include:

\begin{itemize}
\tightlist
\item
  contacting Hannah so we're both aware of this
\end{itemize}

\subsection{20/08/2026:}\label{section-1}

Hannah needs to take a year out; should be fine (diss can be adjusted to
1 person according to Turci), so will go ahead as is).

\subsection{09/09/2026:}\label{section-2}

Have been reading through 100page intro paper on open Qm systems; have
also installed/setup Zotero; its a REALLY nice referencing system, i
think i'm gonna use that for the project.

\subsection{10/09/2026}\label{section-3}

Have received some prereading for tmrs meeting (in person). Its short
but gives bare bone basics.

The bare bones intro states: ``This project will investigate how
different operational definitions of quantum Markovianity appear in a
single impurity model coupled to a structured environment.''

The project sketch outlines first steps:

\subsubsection{part 1 (basic setup, direct
diagonalisation):}\label{part-1-basic-setup-direct-diagonalisation}

\begin{itemize}
\tightlist
\item
  simulate a ``single impurity'' directly
\item
  take a ``fermionic system'' w/single site coupled to a bath. This
  setup involves a single system site coupled to a bunch of bath sites.
  The combined system hamiltonian follows:
\end{itemize}

\$ \hat{H}=\BLU{\epsilon \hat{s}^{\dagger} \hat{s}}+
\ORG{\sum_{k=1}^N g_k\left(\hat{b}_k^{\dagger} \hat{s}+\hat{s}^{\dagger} \hat{b}_k\right)}
+\GRN{\sum_{k=1}^N \omega_k \hat{b}_k^{\dagger} \hat{b}_k}. \$

The first (blue) term is the ``on site'' energy, the second is the
system/bath coupling energy (hopping energy?), with the last term being
the bath's free energy. All these terms are described in terms of the
systems and baths creation/annihalation operators (\(\hat{s}\) and
\(\hat{b}\) respectively). There's also some relation between the
coupling \(g_{k}\) and the ``spectral density''. Also also since only
creation/annhilation operators are involved; can consider the
hamiltonian non interaction (from reading around topic this prob relates
to interaction picture; which is nice that i wont need to deal with that
atm (i think interpreation is that bath sites don't really interact with
each other; and system site (of which there is only 1) doesnt interact
with other system sites (because only 1)).

Fermionic anticommutation relations apply; determine stats.

The ``all bath sites'' connect to the single system site is called a
``star geometry''.

Direct diagonalisation and time evolution evaluation of this system
doesn't scale well w/bath size; so isn't useful for examining continuum
limit (Note; if its quick, may still be worth implementing a solver for
this; firstly as an exercise to familiarise myself w/problem setup, and
to directly illustrate these scaling challenges).

\subsubsection{part 2 (correlation
matrix):}\label{part-2-correlation-matrix}

Given non interacting hamiltonian; correlation matrix of system
determines dynamics. Given an ``operator vector'':

\$
\(\hat{\mathbf{A}} = (\hat{s}, \hat{b}_1, \hat{b}_2, ..., \hat{b}_N)^T\)
\$

Correlation matrix follows:

\(\mathbf{C}_{ij} = \langle \hat{A}_j^\dagger \hat{A}_i \rangle\)

These relate to the hamiltonian via a ``single particle hamiltonian''
\(h\) (unsure exactly what this represents):

\$ \hat{H}=\hat{\mathbf{A}}\^{}\{\dagger\} \mathbf{h}
\hat{\mathbf{A}}=\sum\emph{\{i, j=1\}\^{}\{\infty\} h}\{i j\}
\hat{A}\_i\^{}\{\dagger\} \hat{A}\_j \$

The correlation matrix evolves as:

\(\mathbf{C}(t) = \mathbf{u}(t) \mathbf{C}(0) \mathbf{u}(t)^{\dagger}\)
\(\mathbf{u}(t) = e^{-i \mathbf{h} t}\)

Since the system is what is of interest, and the first term in the
operator vector \(\hat{A}\) represents the system; the reduced density
operator for the system only requires the first term in the correlation
matrix as:

\$ \rho\_S(t)=\left(

\begin{array}{cc}
1-\mathbf{C}_{11}(t) & 0 \\
0 & \mathbf{C}_{11}(t)
\end{array}

\right) \$

I suspect in practice, the starting system state would inform
\(C_{11}=\langle \hat{s}^{}\hat{s} \rangle\), and time evolution on
\(C_{11}\) would be performed using the system single particle
hamiltonian; then the reduced density operator can be computed for any
time and examined. The specifics of how this is to be done are unclear
to me atm.

\subsubsection{part 3 (``chain mapping''):}\label{part-3-chain-mapping}

Chain mapping maps system from ``star geometry'' to ``semi-infinite'' 1D
chain. proceeds as: * discretize bath spectral density to determine bath
couplings/modes * unitary transform maps bath onto 1D chain with nearest
neighbour couplings (seemingly trying to ``match'' modes of complete
bath?) - system couples at start of chain * resulting hamiltonian is
seemingly nicely behaved and easy to simulate? * Note finite length of
chain can cause ``reflections'' which can mess w/signal - there are ways
to deal w/this

\subsubsection{part 4 (extracting dynamical
map):}\label{part-4-extracting-dynamical-map}

Key point of these sims is to extract ``dynamical map'' \(\Gamma\) which
maps the initial reduced density \(\rho_{s}(0)\) to \(\rho_{s}(t)\) as
\(\rho_{s}(t) = \Gamma(t, 0)\rho_{s}(0)\). This dynamical map becomes
main ``result'' of simulations and is necessary for analysis.

\subsubsection{part 5 (testing memory
effects):}\label{part-5-testing-memory-effects}

The dynamical map can then be examined to test for memory effects using
various different metrics: * memory kernels - defined by
``Nakajima-Zwanzig formalism'' * Information backflow - defined by BPL
(Breuer-Laine-Piilo) measure * Failure of Completely positive (CP)
divisibility - RHP (Rivas-Huelga-Plenio) measure

I'm unsure exactly what these different measures capture, but seem to be
main ``probes'' into the system's memory effects.

\subsection{Meeting discussion points:}\label{meeting-discussion-points}

\begin{itemize}
\tightlist
\item
  Physics discussion:

  \begin{itemize}
  \tightlist
  \item
    specifics of simulation setup
  \item
    basics reading (i think paper i have and book i have are these)

    \begin{itemize}
    \tightlist
    \item
      the references in this project sketch may be useful
    \end{itemize}
  \item
    associated research area reading (papers that sketch out research
    area)
  \end{itemize}
\item
  Admin discussion:

  \begin{itemize}
  \tightlist
  \item
    What do i need to setup:

    \begin{itemize}
    \tightlist
    \item
      Ref manager (Zotero) {[}x{]}
    \item
      Julia {[}x{]}
    \item
      IDE (VScodium) {[}x{]}
    \item
      Julia Itensor package {[}x{]}
    \item
      Labbook (default instinct is quarto, but want to ask)
    \end{itemize}
  \item
    future meetings and general timeline

    \begin{itemize}
    \tightlist
    \item
      when regular supervisor meetings
    \item
      when regular time to work on project
    \item
      what are project stages (I.E, general research,
      coding/experimentation..?)
    \item
      what milestones are part of diss (i forget exactly where i saw em,
      but there were presentations\ldots{} that i need to do early)
    \end{itemize}
  \end{itemize}
\end{itemize}

11/09/2026 meeting Recommended Condensed matter field theory chapter 2
(second quantisation)

compare measures for non markoviantity; qualitative difference; edge
cases of where measures fail

programming language of choice; julia or python

general explanation of setup

current plan for recurring meeting is thursday 10:00-11:00 potentially.

To do list: * setup labbook as quarto notebook * continue familiarise
self w/Julia * continue reading paper * experiment w/code Turci will
send me




\end{document}
